% --- Archon physics formalization source begin ---
% archon:physics
% archon:covers PhyXMiniProblems/problem_phyx_mini_0646.lean
% archon:source-report reports/phyx_mini/problem_phyx_mini_0646.source.json

\chapter{Physics problem phyx\_mini\_0646}
\label{ch:PhyXMiniProblems_problem_phyx_mini_0646}

\paragraph{Problem source.}
\#\# Physical scenario

Figure shows the spectrum of light emitted by a firefly.

\#\# Question

Determine the temperature of a blackbody that would emit radiation peaked at the same frequency.

\#\# Answer choices

- A: A: \textbackslash{}( 2\textbackslash{}

- B: B: \textbackslash{}( 1\textbackslash{}

- C: C: \textbackslash{}( 5\textbackslash{}

- D: D: \textbackslash{}( 4\textbackslash{}

\#\# Figure caption (auxiliary; use the image as primary evidence)

The image is a graph showing a plot of relative intensity versus wavelength in nanometers (nm). The x-axis is labeled "Wavelength (nm)" and ranges from around 300 to 700 nm, with ticks at 400, 500, 600, and 700 nm. The y-axis is labeled "Relative intensity" and ranges from 0.0 to 1.0 with ticks at intervals of 0.2. 

A curve on the graph is shaded in a light yellow color and represents the relative intensity across different wavelengths. It rises sharply, peaks slightly above 600 nm, and then falls back down more gradually. The peak intensity reaches 1.0. The outline of the curve is drawn in orange. 

This graph likely represents spectrum data, showing how intensity varies with wavelength within the visible spectrum.

\#\# Dataset metadata

Quantum Phenomena; Physical Model Grounding Reasoning, Multi-Formula Reasoning

\paragraph{Recorded answer/context.}
C: C: \textbackslash{}( 5\textbackslash{}

\paragraph{Figure/image path.}
/root/proposal\_for\_physic/science-mango/phyx\_mini\_run/phyx\_data/test\_image/646.png

\paragraph{Formalization target.}
create a compiling Lean file with sorry bodies at `PhyXMiniProblems/problem\_phyx\_mini\_0646.lean`. The Lean declarations must preserve the physical quantities, dimensions or dimensional roles, figure labels, governing-law hypotheses, and final relation expressed by this problem.
Use Mathlib/Physlib names found through LeanExplore where available. If a physics API is missing, introduce faithful local abstractions rather than scalar placeholder aliases.

\begin{theorem}[Physics formalization target]
\label{thm:physics:phyx_mini_0646:target}
\lean{PhyXMiniProblems.ProblemPhyXMini0646.problem_phyx_mini_0646}
\uses{def:physics:phyx-mini-0646:phyxminiproblems-problemphyxmini0646-wavelengthinmeters, def:physics:phyx-mini-0646:phyxminiproblems-problemphyxmini0646-wienconstantinmeterkelvins, def:physics:phyx-mini-0646:phyxminiproblems-problemphyxmini0646-temperatureinkelvins, def:physics:phyx-mini-0646:phyxminiproblems-problemphyxmini0646-matchesprimaryraster646, def:physics:phyx-mini-0646:phyxminiproblems-problemphyxmini0646-fireflyblackbodypeaksetup, def:physics:phyx-mini-0646:phyxminiproblems-problemphyxmini0646-matchesfireflyblackbodyscenario, def:physics:phyx-mini-0646:phyxminiproblems-problemphyxmini0646-matchesreportedfireflyspectrumcaption, def:physics:phyx-mini-0646:phyxminiproblems-problemphyxmini0646-hasmeasuredfireflyspectrumpeak, def:physics:phyx-mini-0646:phyxminiproblems-problemphyxmini0646-hasphysicalpeakmatchingparameters, def:physics:phyx-mini-0646:phyxminiproblems-problemphyxmini0646-usesstandardblackbodyconstants, def:physics:phyx-mini-0646:phyxminiproblems-problemphyxmini0646-satisfiesvacuumdispersionandwienlaws, def:physics:phyx-mini-0646:phyxminiproblems-problemphyxmini0646-hassamepeakfrequency}
The assigned autoformalize agent should translate this physics problem into Lean declarations in the covered file, with theorem and lemma proof bodies written as `by sorry`.
\end{theorem}
\begin{proof}
This is an autoformalization task, not a proof task. Produce statements that can later be proved by the physics prover without weakening the physical model.
\end{proof}

% --- Archon declaration topology begin ---
\section{Declaration topology}
\begin{definition}[frequency Dimension]
  \label{def:physics:phyx-mini-0646:phyxminiproblems-problemphyxmini0646-frequencydimension}
  \lean{PhyXMiniProblems.ProblemPhyXMini0646.frequencyDimension}
  Frequency has dimension inverse time
\end{definition}
\begin{proof}
  This is the declared model definition of frequency Dimension.
\end{proof}
\begin{definition}[speed Dimension]
  \label{def:physics:phyx-mini-0646:phyxminiproblems-problemphyxmini0646-speeddimension}
  \lean{PhyXMiniProblems.ProblemPhyXMini0646.speedDimension}
  Speed has dimension length per time
\end{definition}
\begin{proof}
  This is the declared model definition of speed Dimension.
\end{proof}
\begin{definition}[probability Density Dimension]
  \label{def:physics:phyx-mini-0646:phyxminiproblems-problemphyxmini0646-probabilitydensitydimension}
  \lean{PhyXMiniProblems.ProblemPhyXMini0646.probabilityDensityDimension}
  A one-dimensional probability density has dimension inverse length
\end{definition}
\begin{proof}
  This is the declared model definition of probability Density Dimension.
\end{proof}
\begin{definition}[wien Displacement Constant Dimension]
  \label{def:physics:phyx-mini-0646:phyxminiproblems-problemphyxmini0646-wiendisplacementconstantdimension}
  \lean{PhyXMiniProblems.ProblemPhyXMini0646.wienDisplacementConstantDimension}
  Wien's wavelength displacement constant has dimension length-temperature
\end{definition}
\begin{proof}
  This is the declared model definition of wien Displacement Constant Dimension.
\end{proof}
\begin{definition}[Wavelength Quantity]
  \label{def:physics:phyx-mini-0646:phyxminiproblems-problemphyxmini0646-wavelengthquantity}
  \lean{PhyXMiniProblems.ProblemPhyXMini0646.WavelengthQuantity}
  A nonnegative physical wavelength
\end{definition}
\begin{proof}
  This is the declared model definition of Wavelength Quantity.
\end{proof}
\begin{definition}[Position Quantity]
  \label{def:physics:phyx-mini-0646:phyxminiproblems-problemphyxmini0646-positionquantity}
  \lean{PhyXMiniProblems.ProblemPhyXMini0646.PositionQuantity}
  A signed physical position, used for the primary raster's x coordinate
\end{definition}
\begin{proof}
  This is the declared model definition of Position Quantity.
\end{proof}
\begin{definition}[Frequency Quantity]
  \label{def:physics:phyx-mini-0646:phyxminiproblems-problemphyxmini0646-frequencyquantity}
  \lean{PhyXMiniProblems.ProblemPhyXMini0646.FrequencyQuantity}
  \uses{def:physics:phyx-mini-0646:phyxminiproblems-problemphyxmini0646-frequencydimension}
  A nonnegative physical frequency
\end{definition}
\begin{proof}
  This is the declared model definition of Frequency Quantity.
\end{proof}
\begin{definition}[Speed Quantity]
  \label{def:physics:phyx-mini-0646:phyxminiproblems-problemphyxmini0646-speedquantity}
  \lean{PhyXMiniProblems.ProblemPhyXMini0646.SpeedQuantity}
  \uses{def:physics:phyx-mini-0646:phyxminiproblems-problemphyxmini0646-speeddimension}
  A nonnegative physical speed
\end{definition}
\begin{proof}
  This is the declared model definition of Speed Quantity.
\end{proof}
\begin{definition}[Probability Density Quantity]
  \label{def:physics:phyx-mini-0646:phyxminiproblems-problemphyxmini0646-probabilitydensityquantity}
  \lean{PhyXMiniProblems.ProblemPhyXMini0646.ProbabilityDensityQuantity}
  \uses{def:physics:phyx-mini-0646:phyxminiproblems-problemphyxmini0646-probabilitydensitydimension}
  A nonnegative one-dimensional probability density
\end{definition}
\begin{proof}
  This is the declared model definition of Probability Density Quantity.
\end{proof}
\begin{definition}[Wien Displacement Constant Quantity]
  \label{def:physics:phyx-mini-0646:phyxminiproblems-problemphyxmini0646-wiendisplacementconstantquantity}
  \lean{PhyXMiniProblems.ProblemPhyXMini0646.WienDisplacementConstantQuantity}
  \uses{def:physics:phyx-mini-0646:phyxminiproblems-problemphyxmini0646-wiendisplacementconstantdimension}
  A physical value of Wien's wavelength displacement constant
\end{definition}
\begin{proof}
  This is the declared model definition of Wien Displacement Constant Quantity.
\end{proof}
\begin{definition}[nnreal Quantity Readout]
  \label{def:physics:phyx-mini-0646:phyxminiproblems-problemphyxmini0646-nnrealquantityreadout}
  \lean{PhyXMiniProblems.ProblemPhyXMini0646.nnrealQuantityReadout}
  Read a nonnegative dimensionful quantity in a coherent unit system
\end{definition}
\begin{proof}
  This is the declared model definition of nnreal Quantity Readout.
\end{proof}
\begin{definition}[real Quantity Readout]
  \label{def:physics:phyx-mini-0646:phyxminiproblems-problemphyxmini0646-realquantityreadout}
  \lean{PhyXMiniProblems.ProblemPhyXMini0646.realQuantityReadout}
  Read a signed dimensionful quantity in a coherent unit system
\end{definition}
\begin{proof}
  This is the declared model definition of real Quantity Readout.
\end{proof}
\begin{definition}[wavelength Readout]
  \label{def:physics:phyx-mini-0646:phyxminiproblems-problemphyxmini0646-wavelengthreadout}
  \lean{PhyXMiniProblems.ProblemPhyXMini0646.wavelengthReadout}
  \uses{def:physics:phyx-mini-0646:phyxminiproblems-problemphyxmini0646-wavelengthquantity, def:physics:phyx-mini-0646:phyxminiproblems-problemphyxmini0646-nnrealquantityreadout}
  Read a wavelength in a selected length unit
\end{definition}
\begin{proof}
  This is the declared model definition of wavelength Readout.
\end{proof}
\begin{definition}[wavelength In Meters]
  \label{def:physics:phyx-mini-0646:phyxminiproblems-problemphyxmini0646-wavelengthinmeters}
  \lean{PhyXMiniProblems.ProblemPhyXMini0646.wavelengthInMeters}
  \uses{def:physics:phyx-mini-0646:phyxminiproblems-problemphyxmini0646-wavelengthquantity, def:physics:phyx-mini-0646:phyxminiproblems-problemphyxmini0646-wavelengthreadout}
  Read a wavelength in metres
\end{definition}
\begin{proof}
  This is the declared model definition of wavelength In Meters.
\end{proof}
\begin{definition}[wavelength In Nanometers]
  \label{def:physics:phyx-mini-0646:phyxminiproblems-problemphyxmini0646-wavelengthinnanometers}
  \lean{PhyXMiniProblems.ProblemPhyXMini0646.wavelengthInNanometers}
  \uses{def:physics:phyx-mini-0646:phyxminiproblems-problemphyxmini0646-wavelengthquantity, def:physics:phyx-mini-0646:phyxminiproblems-problemphyxmini0646-wavelengthreadout}
  Read a wavelength in nanometres
\end{definition}
\begin{proof}
  This is the declared model definition of wavelength In Nanometers.
\end{proof}
\begin{definition}[position In Femtometers]
  \label{def:physics:phyx-mini-0646:phyxminiproblems-problemphyxmini0646-positioninfemtometers}
  \lean{PhyXMiniProblems.ProblemPhyXMini0646.positionInFemtometers}
  \uses{def:physics:phyx-mini-0646:phyxminiproblems-problemphyxmini0646-positionquantity, def:physics:phyx-mini-0646:phyxminiproblems-problemphyxmini0646-realquantityreadout}
  Read a signed position in femtometres
\end{definition}
\begin{proof}
  This is the declared model definition of position In Femtometers.
\end{proof}
\begin{definition}[frequency In Hertz]
  \label{def:physics:phyx-mini-0646:phyxminiproblems-problemphyxmini0646-frequencyinhertz}
  \lean{PhyXMiniProblems.ProblemPhyXMini0646.frequencyInHertz}
  \uses{def:physics:phyx-mini-0646:phyxminiproblems-problemphyxmini0646-frequencyquantity, def:physics:phyx-mini-0646:phyxminiproblems-problemphyxmini0646-nnrealquantityreadout}
  Read a frequency in coherent-SI hertz
\end{definition}
\begin{proof}
  This is the declared model definition of frequency In Hertz.
\end{proof}
\begin{definition}[speed In Meters Per Second]
  \label{def:physics:phyx-mini-0646:phyxminiproblems-problemphyxmini0646-speedinmeterspersecond}
  \lean{PhyXMiniProblems.ProblemPhyXMini0646.speedInMetersPerSecond}
  \uses{def:physics:phyx-mini-0646:phyxminiproblems-problemphyxmini0646-speedquantity, def:physics:phyx-mini-0646:phyxminiproblems-problemphyxmini0646-nnrealquantityreadout}
  Read a speed in coherent-SI metres per second
\end{definition}
\begin{proof}
  This is the declared model definition of speed In Meters Per Second.
\end{proof}
\begin{definition}[probability Density In Per Femtometer]
  \label{def:physics:phyx-mini-0646:phyxminiproblems-problemphyxmini0646-probabilitydensityinperfemtometer}
  \lean{PhyXMiniProblems.ProblemPhyXMini0646.probabilityDensityInPerFemtometer}
  \uses{def:physics:phyx-mini-0646:phyxminiproblems-problemphyxmini0646-probabilitydensityquantity, def:physics:phyx-mini-0646:phyxminiproblems-problemphyxmini0646-nnrealquantityreadout}
  Read a probability density in inverse femtometres
\end{definition}
\begin{proof}
  This is the declared model definition of probability Density In Per Femtometer.
\end{proof}
\begin{definition}[wien Constant In Meter Kelvins]
  \label{def:physics:phyx-mini-0646:phyxminiproblems-problemphyxmini0646-wienconstantinmeterkelvins}
  \lean{PhyXMiniProblems.ProblemPhyXMini0646.wienConstantInMeterKelvins}
  \uses{def:physics:phyx-mini-0646:phyxminiproblems-problemphyxmini0646-wiendisplacementconstantquantity, def:physics:phyx-mini-0646:phyxminiproblems-problemphyxmini0646-nnrealquantityreadout}
  Read Wien's constant in coherent-SI metre-kelvins
\end{definition}
\begin{proof}
  This is the declared model definition of wien Constant In Meter Kelvins.
\end{proof}
\begin{definition}[temperature In Kelvins]
  \label{def:physics:phyx-mini-0646:phyxminiproblems-problemphyxmini0646-temperatureinkelvins}
  \lean{PhyXMiniProblems.ProblemPhyXMini0646.temperatureInKelvins}
  Read a Physlib absolute temperature in kelvins. Temperature stores a nonnegative magnitude in an arbitrary zero-preserving temperature unit, so the storage unit is retained explicitly
\end{definition}
\begin{proof}
  This is the declared model definition of temperature In Kelvins.
\end{proof}
\begin{definition}[Primary Raster Axis]
  \label{def:physics:phyx-mini-0646:phyxminiproblems-problemphyxmini0646-primaryrasteraxis}
  \lean{PhyXMiniProblems.ProblemPhyXMini0646.PrimaryRasterAxis}
  The two axes visible in the primary bitmap
\end{definition}
\begin{proof}
  This is the declared model definition of Primary Raster Axis.
\end{proof}
\begin{definition}[Primary Raster Axis Role]
  \label{def:physics:phyx-mini-0646:phyxminiproblems-problemphyxmini0646-primaryrasteraxisrole}
  \lean{PhyXMiniProblems.ProblemPhyXMini0646.PrimaryRasterAxisRole}
  \uses{def:physics:phyx-mini-0646:phyxminiproblems-problemphyxmini0646-positioninfemtometers}
  Physical roles of the two primary-raster axes
\end{definition}
\begin{proof}
  This is the declared model definition of Primary Raster Axis Role.
\end{proof}
\begin{definition}[Primary Raster Curve Geometry]
  \label{def:physics:phyx-mini-0646:phyxminiproblems-problemphyxmini0646-primaryrastercurvegeometry}
  \lean{PhyXMiniProblems.ProblemPhyXMini0646.PrimaryRasterCurveGeometry}
  Qualitative shape of the green curve visible in image 646.png
\end{definition}
\begin{proof}
  This is the declared model definition of Primary Raster Curve Geometry.
\end{proof}
\begin{definition}[Primary Raster646]
  \label{def:physics:phyx-mini-0646:phyxminiproblems-problemphyxmini0646-primaryraster646}
  \lean{PhyXMiniProblems.ProblemPhyXMini0646.PrimaryRaster646}
  \uses{def:physics:phyx-mini-0646:phyxminiproblems-problemphyxmini0646-positionquantity, def:physics:phyx-mini-0646:phyxminiproblems-problemphyxmini0646-probabilitydensityquantity, def:physics:phyx-mini-0646:phyxminiproblems-problemphyxmini0646-primaryrasteraxis, def:physics:phyx-mini-0646:phyxminiproblems-problemphyxmini0646-primaryrasteraxisrole, def:physics:phyx-mini-0646:phyxminiproblems-problemphyxmini0646-primaryrastercurvegeometry}
  Typed transcription of the actual primary bitmap. The scale a is kept as a physical inverse-length quantity, rather than as an untyped scalar
\end{definition}
\begin{proof}
  This is the declared model definition of Primary Raster646.
\end{proof}
\begin{definition}[Matches Primary Raster646]
  \label{def:physics:phyx-mini-0646:phyxminiproblems-problemphyxmini0646-matchesprimaryraster646}
  \lean{PhyXMiniProblems.ProblemPhyXMini0646.MatchesPrimaryRaster646}
  \uses{def:physics:phyx-mini-0646:phyxminiproblems-problemphyxmini0646-positioninfemtometers, def:physics:phyx-mini-0646:phyxminiproblems-problemphyxmini0646-probabilitydensityinperfemtometer, def:physics:phyx-mini-0646:phyxminiproblems-problemphyxmini0646-primaryraster646}
  Labels, coordinates, and incidence data read directly from the primary raster. This predicate makes no claim that the bitmap is a firefly spectrum
\end{definition}
\begin{proof}
  This is the declared model definition of Matches Primary Raster646.
\end{proof}
\begin{definition}[Emitter Kind]
  \label{def:physics:phyx-mini-0646:phyxminiproblems-problemphyxmini0646-emitterkind}
  \lean{PhyXMiniProblems.ProblemPhyXMini0646.EmitterKind}
  The observed source and the ideal comparison emitter
\end{definition}
\begin{proof}
  This is the declared model definition of Emitter Kind.
\end{proof}
\begin{definition}[Spectrum Axis]
  \label{def:physics:phyx-mini-0646:phyxminiproblems-problemphyxmini0646-spectrumaxis}
  \lean{PhyXMiniProblems.ProblemPhyXMini0646.SpectrumAxis}
  The axes of the firefly spectrum described by the chapter's caption
\end{definition}
\begin{proof}
  This is the declared model definition of Spectrum Axis.
\end{proof}
\begin{definition}[Spectrum Axis Role]
  \label{def:physics:phyx-mini-0646:phyxminiproblems-problemphyxmini0646-spectrumaxisrole}
  \lean{PhyXMiniProblems.ProblemPhyXMini0646.SpectrumAxisRole}
  \uses{def:physics:phyx-mini-0646:phyxminiproblems-problemphyxmini0646-wavelengthinnanometers}
  Physical or dimensionless role assigned to a reported spectrum axis
\end{definition}
\begin{proof}
  This is the declared model definition of Spectrum Axis Role.
\end{proof}
\begin{definition}[Reported Firefly Spectrum Caption]
  \label{def:physics:phyx-mini-0646:phyxminiproblems-problemphyxmini0646-reportedfireflyspectrumcaption}
  \lean{PhyXMiniProblems.ProblemPhyXMini0646.ReportedFireflySpectrumCaption}
  \uses{def:physics:phyx-mini-0646:phyxminiproblems-problemphyxmini0646-spectrumaxis, def:physics:phyx-mini-0646:phyxminiproblems-problemphyxmini0646-spectrumaxisrole}
  Plot metadata asserted by the auxiliary caption. It is intentionally not identified with PrimaryRaster646, because the actual bitmap has different labels and physical content
\end{definition}
\begin{proof}
  This is the declared model definition of Reported Firefly Spectrum Caption.
\end{proof}
\begin{definition}[Firefly Blackbody Peak Setup]
  \label{def:physics:phyx-mini-0646:phyxminiproblems-problemphyxmini0646-fireflyblackbodypeaksetup}
  \lean{PhyXMiniProblems.ProblemPhyXMini0646.FireflyBlackbodyPeakSetup}
  \uses{def:physics:phyx-mini-0646:phyxminiproblems-problemphyxmini0646-wavelengthquantity, def:physics:phyx-mini-0646:phyxminiproblems-problemphyxmini0646-frequencyquantity, def:physics:phyx-mini-0646:phyxminiproblems-problemphyxmini0646-speedquantity, def:physics:phyx-mini-0646:phyxminiproblems-problemphyxmini0646-wiendisplacementconstantquantity, def:physics:phyx-mini-0646:phyxminiproblems-problemphyxmini0646-primaryraster646, def:physics:phyx-mini-0646:phyxminiproblems-problemphyxmini0646-emitterkind, def:physics:phyx-mini-0646:phyxminiproblems-problemphyxmini0646-reportedfireflyspectrumcaption}
  The requested blackbody temperature is an independent physical field. It is not defined from 5000, an answer label, the reported plot, or Wien's law
\end{definition}
\begin{proof}
  This is the declared model definition of Firefly Blackbody Peak Setup.
\end{proof}
\begin{definition}[Matches Firefly Blackbody Scenario]
  \label{def:physics:phyx-mini-0646:phyxminiproblems-problemphyxmini0646-matchesfireflyblackbodyscenario}
  \lean{PhyXMiniProblems.ProblemPhyXMini0646.MatchesFireflyBlackbodyScenario}
  \uses{def:physics:phyx-mini-0646:phyxminiproblems-problemphyxmini0646-fireflyblackbodypeaksetup}
  Qualitative emitter roles stated in the problem
\end{definition}
\begin{proof}
  This is the declared model definition of Matches Firefly Blackbody Scenario.
\end{proof}
\begin{definition}[Matches Reported Firefly Spectrum Caption]
  \label{def:physics:phyx-mini-0646:phyxminiproblems-problemphyxmini0646-matchesreportedfireflyspectrumcaption}
  \lean{PhyXMiniProblems.ProblemPhyXMini0646.MatchesReportedFireflySpectrumCaption}
  \uses{def:physics:phyx-mini-0646:phyxminiproblems-problemphyxmini0646-wavelengthinnanometers, def:physics:phyx-mini-0646:phyxminiproblems-problemphyxmini0646-fireflyblackbodypeaksetup}
  Axis labels, ticks, colors, and approximate peak data asserted by the chapter caption. The peak is placed only between the stated 600 nm tick and the displayed 700 nm endpoint. No sharper tolerance is invented from the word “slightly.”
\end{definition}
\begin{proof}
  This is the declared model definition of Matches Reported Firefly Spectrum Caption.
\end{proof}
\begin{definition}[Has Measured Firefly Spectrum Peak]
  \label{def:physics:phyx-mini-0646:phyxminiproblems-problemphyxmini0646-hasmeasuredfireflyspectrumpeak}
  \lean{PhyXMiniProblems.ProblemPhyXMini0646.HasMeasuredFireflySpectrumPeak}
  \uses{def:physics:phyx-mini-0646:phyxminiproblems-problemphyxmini0646-fireflyblackbodypeaksetup}
  The measured firefly wavelength is a global maximum of the dimensionless relative-intensity spectrum. Mathlib's IsMaxOn supplies the exact maximum predicate
\end{definition}
\begin{proof}
  This is the declared model definition of Has Measured Firefly Spectrum Peak.
\end{proof}
\begin{definition}[Has Physical Peak Matching Parameters]
  \label{def:physics:phyx-mini-0646:phyxminiproblems-problemphyxmini0646-hasphysicalpeakmatchingparameters}
  \lean{PhyXMiniProblems.ProblemPhyXMini0646.HasPhysicalPeakMatchingParameters}
  \uses{def:physics:phyx-mini-0646:phyxminiproblems-problemphyxmini0646-wavelengthinmeters, def:physics:phyx-mini-0646:phyxminiproblems-problemphyxmini0646-frequencyinhertz, def:physics:phyx-mini-0646:phyxminiproblems-problemphyxmini0646-speedinmeterspersecond, def:physics:phyx-mini-0646:phyxminiproblems-problemphyxmini0646-wienconstantinmeterkelvins, def:physics:phyx-mini-0646:phyxminiproblems-problemphyxmini0646-temperatureinkelvins, def:physics:phyx-mini-0646:phyxminiproblems-problemphyxmini0646-fireflyblackbodypeaksetup}
  Positivity and nondegeneracy conditions for the physical quantities
\end{definition}
\begin{proof}
  This is the declared model definition of Has Physical Peak Matching Parameters.
\end{proof}
\begin{definition}[Uses Standard Blackbody Constants]
  \label{def:physics:phyx-mini-0646:phyxminiproblems-problemphyxmini0646-usesstandardblackbodyconstants}
  \lean{PhyXMiniProblems.ProblemPhyXMini0646.UsesStandardBlackbodyConstants}
  \uses{def:physics:phyx-mini-0646:phyxminiproblems-problemphyxmini0646-speedinmeterspersecond, def:physics:phyx-mini-0646:phyxminiproblems-problemphyxmini0646-wienconstantinmeterkelvins, def:physics:phyx-mini-0646:phyxminiproblems-problemphyxmini0646-fireflyblackbodypeaksetup}
  Standard coherent-SI calibrations. They are universal physical constants, not fitted values of the requested temperature
\end{definition}
\begin{proof}
  This is the declared model definition of Uses Standard Blackbody Constants.
\end{proof}
\begin{definition}[Satisfies Vacuum Dispersion And Wien Laws]
  \label{def:physics:phyx-mini-0646:phyxminiproblems-problemphyxmini0646-satisfiesvacuumdispersionandwienlaws}
  \lean{PhyXMiniProblems.ProblemPhyXMini0646.SatisfiesVacuumDispersionAndWienLaws}
  \uses{def:physics:phyx-mini-0646:phyxminiproblems-problemphyxmini0646-wavelengthinmeters, def:physics:phyx-mini-0646:phyxminiproblems-problemphyxmini0646-frequencyinhertz, def:physics:phyx-mini-0646:phyxminiproblems-problemphyxmini0646-speedinmeterspersecond, def:physics:phyx-mini-0646:phyxminiproblems-problemphyxmini0646-wienconstantinmeterkelvins, def:physics:phyx-mini-0646:phyxminiproblems-problemphyxmini0646-temperatureinkelvins, def:physics:phyx-mini-0646:phyxminiproblems-problemphyxmini0646-fireflyblackbodypeaksetup}
  Vacuum dispersion c = λν for both emission peaks and Wien's wavelength displacement law λ\_max T = b for the ideal blackbody. No rounded temperature conclusion occurs in these governing relations
\end{definition}
\begin{proof}
  This is the declared model definition of Satisfies Vacuum Dispersion And Wien Laws.
\end{proof}
\begin{definition}[Has Same Peak Frequency]
  \label{def:physics:phyx-mini-0646:phyxminiproblems-problemphyxmini0646-hassamepeakfrequency}
  \lean{PhyXMiniProblems.ProblemPhyXMini0646.HasSamePeakFrequency}
  \uses{def:physics:phyx-mini-0646:phyxminiproblems-problemphyxmini0646-fireflyblackbodypeaksetup}
  The design condition from the question: the ideal blackbody is chosen to peak at the same physical frequency as the firefly
\end{definition}
\begin{proof}
  This is the declared model definition of Has Same Peak Frequency.
\end{proof}
\begin{lemma}[matched peak wavelengths equal]
  \label{lem:physics:phyx-mini-0646:phyxminiproblems-problemphyxmini0646-matched-peak-wavelengths-equal}
  \lean{PhyXMiniProblems.ProblemPhyXMini0646.matched_peak_wavelengths_equal}
  \uses{def:physics:phyx-mini-0646:phyxminiproblems-problemphyxmini0646-wavelengthinmeters, def:physics:phyx-mini-0646:phyxminiproblems-problemphyxmini0646-fireflyblackbodypeaksetup, def:physics:phyx-mini-0646:phyxminiproblems-problemphyxmini0646-hasphysicalpeakmatchingparameters, def:physics:phyx-mini-0646:phyxminiproblems-problemphyxmini0646-satisfiesvacuumdispersionandwienlaws, def:physics:phyx-mini-0646:phyxminiproblems-problemphyxmini0646-hassamepeakfrequency}
  Equal positive peak frequencies under c = λν give equal wavelengths
\end{lemma}
\begin{proof}
  The conclusion follows from the governing relations and figure readouts named in the statement of matched peak wavelengths equal.
\end{proof}
\begin{lemma}[matched peak wien relation]
  \label{lem:physics:phyx-mini-0646:phyxminiproblems-problemphyxmini0646-matched-peak-wien-relation}
  \lean{PhyXMiniProblems.ProblemPhyXMini0646.matched_peak_wien_relation}
  \uses{def:physics:phyx-mini-0646:phyxminiproblems-problemphyxmini0646-wavelengthinmeters, def:physics:phyx-mini-0646:phyxminiproblems-problemphyxmini0646-wienconstantinmeterkelvins, def:physics:phyx-mini-0646:phyxminiproblems-problemphyxmini0646-temperatureinkelvins, def:physics:phyx-mini-0646:phyxminiproblems-problemphyxmini0646-fireflyblackbodypeaksetup, def:physics:phyx-mini-0646:phyxminiproblems-problemphyxmini0646-hasphysicalpeakmatchingparameters, def:physics:phyx-mini-0646:phyxminiproblems-problemphyxmini0646-satisfiesvacuumdispersionandwienlaws, def:physics:phyx-mini-0646:phyxminiproblems-problemphyxmini0646-hassamepeakfrequency}
  The exact, unrounded temperature relation obtained after eliminating the blackbody peak wavelength
\end{lemma}
\begin{proof}
  The conclusion follows from the governing relations and figure readouts named in the statement of matched peak wien relation.
\end{proof}
\begin{definition}[Answer Choice]
  \label{def:physics:phyx-mini-0646:phyxminiproblems-problemphyxmini0646-answerchoice}
  \lean{PhyXMiniProblems.ProblemPhyXMini0646.AnswerChoice}
  Labels retained from the truncated multiple-choice source
\end{definition}
\begin{proof}
  This is the declared model definition of Answer Choice.
\end{proof}
\begin{definition}[recorded Answer Label]
  \label{def:physics:phyx-mini-0646:phyxminiproblems-problemphyxmini0646-recordedanswerlabel}
  \lean{PhyXMiniProblems.ProblemPhyXMini0646.recordedAnswerLabel}
  \uses{def:physics:phyx-mini-0646:phyxminiproblems-problemphyxmini0646-answerchoice}
  The source metadata records label C; this value is not used as a premise
\end{definition}
\begin{proof}
  This is the declared model definition of recorded Answer Label.
\end{proof}
% --- Archon declaration topology end ---
% --- Archon physics formalization source end ---
